% ==============================================================================
% Plantilla Institucional de Trabajo Terminal — UPIIZ IPN
% Archivo: ejemplos/codigo.tex
% ==============================================================================
% Ejemplos prácticos para listados de código fuente en C, C++, Python, MATLAB y Arduino.
% ==============================================================================

\section{Guía de ejemplos para inserción de código fuente}
\label{sec:ejemplo-codigo}

% ------------------------------------------------------------------------------
% 1. CÓDIGO EN C / C++ / ARDUINO
% ------------------------------------------------------------------------------
En el \lstlistingname~\ref{lst:codigo-pid-c} se muestra la implementación discreta de una rutina de control PID para microcontroladores.

\begin{lstlisting}[language=C, caption={Rutina de control PID discreto en C/C++.}, label={lst:codigo-pid-c}]
/* Controlador PID discreto con anti-windup */
float calcularPID(float ref, float med, float dt) {
    static float error_prev = 0.0f;
    static float integral = 0.0f;
    
    float error = ref - med;
    integral += error * dt;
    float derivada = (error - error_prev) / dt;
    error_prev = error;
    
    float salida = (KP * error) + (KI * integral) + (KD * derivada);
    
    // Saturacion de la senal de control
    if (salida > PWM_MAX) salida = PWM_MAX;
    else if (salida < PWM_MIN) salida = PWM_MIN;
    
    return salida;
}
\end{lstlisting}

% ------------------------------------------------------------------------------
% 2. CÓDIGO EN PYTHON
% ------------------------------------------------------------------------------
En el \lstlistingname~\ref{lst:codigo-python} se ilustra un script en Python para filtrado de datos en tiempo real.

\begin{lstlisting}[language=Python, caption={Filtrado digital pasa-bajas recursivo en Python.}, label={lst:codigo-python}]
import numpy as np

def filtro_pasabajas(senal_cruda, alfa=0.15):
    """Aplica un filtro IIR exponencial de primer orden."""
    filtrada = np.zeros_like(senal_cruda)
    filtrada[0] = senal_cruda[0]
    for i in range(1, len(senal_cruda)):
        filtrada[i] = alfa * senal_cruda[i] + (1.0 - alfa) * filtrada[i-1]
    return filtrada
\end{lstlisting}

% ------------------------------------------------------------------------------
% 3. CÓDIGO EN MATLAB
% ------------------------------------------------------------------------------
En el \lstlistingname~\ref{lst:codigo-matlab} se presenta un script de simulación de respuesta al escalón en MATLAB.

\begin{lstlisting}[language=Matlab, caption={Simulación y sintonización de parámetros en MATLAB.}, label={lst:codigo-matlab}]
% Definicion de la planta mecatronica
m = 3.5;  % Masa en kg
b = 1.2;  % Coeficiente de amortiguamiento
k = 25.0; % Constante elastica

num = [1];
den = [m, b, k];
G = tf(num, den);

% Sintonia del controlador PID
C = pid(120, 35, 12);
LazoCerrado = feedback(C*G, 1);
step(LazoCerrado);
grid on;
title('Respuesta al Escalon del Sistema');
\end{lstlisting}
